\documentclass[11pt]{article}
\usepackage[margin=1in]{geometry}
\usepackage{amsmath,amssymb,parskip}
\usepackage[colorlinks=true,allcolors=blue]{hyperref}
\newcommand{\E}{\mathbb E}
\title{Event-Aware Equity-Option Pricing}
\author{}\date{September 2026}
\begin{document}
\maketitle

\section{Structural model}
Let $S_0$ be spot, $K$ strike, $T$ elapsed calendar years to expiry, $r$ the rate,
and $q$ dividend yield. Write $F=S_0e^{(r-q)T}$ and $D=e^{-rT}$.
A scheduled event affects an option only when its event time lies after valuation
and no later than expiry. For flat diffusion and independent scheduled jumps,
\[
 \log(S_T/S_0)=(r-q)T-\tfrac12\sigma_d^2\tau
       +\sigma_d\sqrt{\tau}Z+\sum_j J_j.
\]
Here $\tau$ is variance time. Calendar time governs carry and discounting, while
$\tau$ may weight exchange sessions, overnight closures, weekends and holidays
differently. The default calendar clock has $\tau=T$. Clock weights inferred
from historical stock returns are physical estimates, not automatically
risk-neutral calibration parameters.

Each jump can have a Gaussian-mixture pricing distribution,
\[
 f_J(x)=\sum_{i=1}^{M}w_i\mathcal N(x;\mu_i,s_i^2),\qquad
 w_i\ge0,\quad\sum_iw_i=1.
\]
Given raw centres $m_i$, the common shift
\[
 \mu_i=m_i-\log\left(\sum_iw_i e^{m_i+s_i^2/2}\right)
\]
ensures $\E^Q[e^J]=1$. This is a pricing-measure martingale condition. The
components are not identified as EPS beat/miss outcomes, and their weights are
not estimates of physical earnings probabilities.

\section{Closed-form mixture prices}
With one event and flat diffusion, condition on the mixture component. The
terminal log return is Gaussian on each component, so
\[
 C=D\sum_iw_i\{F_i\Phi(d_{1i})-K\Phi(d_{2i})\},
\]
where
\[
 F_i=Fe^{\mu_i+s_i^2/2},\quad B_i=\sigma_d^2\tau+s_i^2,\quad
 d_{1i}=\frac{\log(F_i/K)+B_i/2}{\sqrt{B_i}},\quad d_{2i}=d_{1i}-\sqrt{B_i}.
\]
Put prices follow from parity. Conditional Gaussian structure makes this finite
mixture formula possible; jumps in general do not necessarily require Monte
Carlo, nor does a fixed event date alone guarantee such a formula.

\section{Heston and Fourier pricing}
The optional Heston variance process evolves in variance time,
\[
 dv_\tau=\kappa(\theta-v_\tau)d\tau+\xi\sqrt{v_\tau}\,dW^v_\tau,
 \qquad d\langle W^S,W^v\rangle_\tau=\rho\,d\tau.
\]
Under the assumed independence of scheduled jumps and the diffusion,
\[
 \varphi_{\rm total}(u)=\varphi_{\rm Heston}(u;\tau)
     e^{iu(r-q)T}\prod_j\varphi_{J_j}(u),\qquad
 \varphi_J(u)=\sum_iw_i e^{iu\mu_i-u^2s_i^2/2},
\]
when the Heston characteristic function is evaluated with zero carry. The COS
method recovers option values from this characteristic function. Multiplying
jump characteristic functions costs $O(N\sum_j M_j)$ component evaluations at
$N$ Fourier frequencies, rather than enumerating $\prod_j M_j$ joint scenarios.
That statement concerns the characteristic-function calculation, not the entire
calibration runtime. Numerical truncation still requires care.

The research estimator fixes $\theta=v_0=\sigma^2$ and $\kappa=2$, and fits
$\sigma,\xi,\rho$ and an event scale with a prescribed mixture shape. It is a
restricted candidate specification, not a universal Heston calibration.
Some studies fit only snapshots with an expiry before and after the next event;
this is a conservative identification rule, not proof that other term structures
contain no event information. Later experiments carry an earlier identifiable
estimate forward or fit a simpler variance term structure.

\section{Market-anchored pricing}
Reference contracts determine forwards and observed implied volatilities.
Target strikes and their parity twins are withheld. For the structural hybrid,
\[
 \sigma_{\rm prior}(k)=\sigma_{\rm structural}(k)
 +\operatorname{interp}\bigl[\sigma_{\rm reference}
                           -\sigma_{\rm structural,reference}\bigr](k).
\]
The market-only alternative uses PCHIP interpolation of reference IV. These
priors are converted to normalized call prices $c(x)=C/(DF)$, $x=K/F$.
A natural cubic spline is fitted on reference strikes and intervening midpoints,
penalizing reference-price error and deviation from the prior between anchors.
Reference errors are scaled by half-spreads with a small floor.

Linear derivative restrictions enforce convexity and decreasing call prices:
$c''\ge0$ and $-1\le c'\le0$, together with payoff bounds. Explicit convex tails
continue the curve outside its anchor range. Put prices use the same forward
and parity. These are within-expiry restrictions; the construction does not
claim absence of calendar arbitrage across all maturities, nor does it model
American early exercise.

\section{Event-preserving repricing}
A simple event model writes total implied variance approximately as
\[
 w(T)=aT+b\,n(T),\qquad a,b\ge0,
\]
where $n(T)$ counts included events. An ATM fit, a two-expiry difference, or a
strike-dependent fit gives alternative estimates. A structural estimate can
instead use the implied-variance difference between model prices with and
without an event, rather than equating log-jump variance with implied variance.

A transport experiment refreshes the current reference smile, keeps its
background variance rate, and removes the estimated event term after the
release. Every method receives the later stock price at scoring time. This is
conditional option repricing, not a stock-return forecast. With only release
dates, a fixed two-session bracket covers the possible announcement timing;
it is never chosen using the larger realized return.

\section{Empirical outcome and scope}
The final structural hybrid has mean absolute error of 4.849 spot basis points
on 23,650 withheld quotes from 251 eligible snapshots in January--August 2025;
87.0\% lie inside bid--ask. The market-only constrained mode scores 4.811 bp and
89.7\%. These methods are competitive with market interpolation, not a
consistent improvement over it. Complete tables, strategy outcomes and sample
limits are in \texttt{PROJECT.md} and \texttt{final\_metrics.json}.

No robust after-cost trading advantage or superior structural hedge was
established. Pricing-measure quantities do not identify physical fair values;
a market-versus-model residual is not automatically mispricing. The ingredients
have prior literature. Scheduled earnings jumps, Heston, Fourier pricing,
variance clocks and constrained interpolation are not claimed as new inventions.
Historical event schedules, European exercise assumptions and daily execution
scenarios limit the scope of the empirical claims.

\end{document}
